% ===================================================================
%  TAULA N.º 3 — Infància i Joventut.
%  Traduction 2026 d'apres texto/io, controlee sur texto/fr et
%  texto/es. Consigne : voir texto/ca/10-taula-01.tex.
% ===================================================================

%%K t03-apar-1 apar
\VUsaut{3.0mm}
\VUtitre{15.0pt}{60}{TAULA N.º 3}
\VUsaut{1.6mm}
\VUtitre{11.4pt}{0}{Infància i Joventut.}
\VUsaut{3.4mm}

%%K t03-apar-2 apar
(Les taules 3 i 4 estan disposades de manera que presenten, en quatre
\VUgras{escenes familiars}, el vocabulari essencial del \VUgras{temps}, l'\VUgras{edat},
la \VUgras{família}, la \VUgras{roba} i la \VUgras{salut}.)
\VUsaut{4.0mm}

%%K t03-tit-1 sub
\VUcentre{12.0pt}{0}{\textit{Primera escena.}}
\VUsaut{3.0mm}

%%K t03-tit-2 sub
\VUpk{10.2pt}{40}{Un bateig al poble.}
\VUsaut{2.4mm}

%%K t03-tit-3 sub
\VUpk{10.2pt}{40}{La primavera.}
\VUsaut{3.0mm}

%%K t03-c1-01-1 p
1. --- Som a la \VUgras{primavera}, al mes de \VUgras{març}, d'\VUgras{abril} o de
\VUgras{maig}, en un dia de la \VUgras{setmana} --- \VUgras{dilluns},
\VUgras{dimarts} o \VUgras{dimecres}. Són les \VUgras{deu} del matí.

%%K t03-c1-02-1 p
2. --- Fa bon \VUgras{temps}, l'\VUgras{aire} és pur. El sol brilla. Uns
\VUgras{nuvolets} \textsuperscript{(2)} pugen pel \VUgras{cel}
\textsuperscript{(1)} blau.

%%K t03-c1-03-1 p
3. --- Els \VUgras{arbres} amb prou feines tenen \VUgras{fulles}. Els prims
\VUgras{pollancres} \textsuperscript{(3)} i l'alt \VUgras{plàtan}
\textsuperscript{(17)} encara no porten més que borrons.

%%K t03-c1-04-1 p
4. --- Les \VUgras{orenetes} \textsuperscript{(4)} han tornat i voleien amb alegres xiscles al
voltant de la \VUgras{fletxa} \textsuperscript{(7)} del
\VUgras{campanar} \textsuperscript{(5)} que corona l'\VUgras{església}
\textsuperscript{(6)} del poble.

%%K t03-tit-4 sub
\VUsaut{3.4mm}
\VUpk{10.2pt}{40}{L'Església.}
\VUsaut{3.0mm}

%%K t03-c1-05-1 p
5. --- Molt senzilla és aquesta església amb el seu gran \VUgras{portal}
\textsuperscript{(13)} i el seu gruixut \VUgras{rellotge} \textsuperscript{(8)}.
L'\VUgras{agulla petita} \textsuperscript{(9)} és sobre la xifra X: marca les
\VUgras{hores}. La \VUgras{gran} \textsuperscript{(10)} és sobre les XII
(migdia); marca els \VUgras{minuts}.

%%K t03-c1-06-1 p
6. --- Al campanar, les \VUgras{campanes} \textsuperscript{(11)} sonen alegrement. Dalt de la
teulada s'alça una petita \VUgras{creu} \textsuperscript{(12)} de pedra.

%%K t03-c1-07-1 p
7. --- L'església no té només un gran \VUgras{portal} \textsuperscript{(13)}, té també una
porteta lateral. Per aquesta porta el senyor \VUgras{rector} \textsuperscript{(15)}, vestit
amb la seva \VUgras{sotana}, surt de la
\VUgras{sagristia} \textsuperscript{(14)} i es dirigeix a la \VUgras{rectoria}
\textsuperscript{(19)}.

%%K t03-c1-08-1 p
8. --- Al seu costat hi ha la \VUgras{lletera} \textsuperscript{(24)} amb el seu
\VUgras{ase} \textsuperscript{(23)}. Torna de la ciutat, on ha portat la seva bona
llet per als nens.

%%K t03-tit-5 sub
\VUsaut{4.0mm}
\VUpk{10.2pt}{40}{L'hort de fruiters.}
\VUsaut{3.4mm}

%%K t03-c1-09-1 p
9. --- El senyor Aliboró (l'ase) està molt content d'aquesta correguda matinal. Aspira amb
delit l'aire embalsamat pel perfum de les \VUgras{flors} \textsuperscript{(20)} que cobreixen
els \VUgras{arbres fruiters} \textsuperscript{(18)} de l'\VUgras{hort} veí. Tots rosats i
blancs són aquells
\VUgras{presseguers} \textsuperscript{(18)} i aquells \VUgras{albercoquers}
\textsuperscript{(19)}, que prometen una bona collita.

%%K t03-c1-10-1 p
10. --- Mentrestant, les petites \VUgras{abelles} \textsuperscript{(22)} del
\VUgras{rusc} \textsuperscript{(21)} hi van a libar la seva dolça \VUgras{mel},
brunzint.

%%K t03-c1-11-1 p
11. --- Però què passa? Són els nens del poble, que arriben corrents i fan fugir les
\VUgras{oques} \textsuperscript{(25)} espantades. És el seguici que surt de l'església després
del \VUgras{bateig} del fill del notari.

%%K t03-c1-12-1 p
12. --- El petit Jaume, al seu \VUgras{bressol} \textsuperscript{(26)}, es desperta amb el so
alegre de les campanes, i la seva germana Magdalena, que també vol veure el seguici, demana a
la seva mare que la vingui a pentinar i a vestir al llindar de la casa.

%%K t03-c1-13-1 p
13. --- La mare hi accedeix. Es posa el seu \VUgras{mocador de cap} \textsuperscript{(28)}, el
seu \VUgras{cosset} \textsuperscript{(29)} i el seu
\VUgras{davantal} \textsuperscript{(30)}, i ve a seure en una cadireta davant la
porta. Magdalena només porta els seus \VUgras{enagos} \textsuperscript{(31)} i el seu
\VUgras{cotilla} \textsuperscript{(32)}.

%%K t03-c1-14-1 p
14. --- Que feliç que és! S'oblida de les seves flors preferides, els seus
\VUgras{clavells} \textsuperscript{(34)}, on es posen les \VUgras{papallones}
\textsuperscript{(35)}, els seus \VUgras{tulipans} \textsuperscript{(36)}, els seus
\VUgras{muguets} \textsuperscript{(37)} en \VUgras{testos} \textsuperscript{(38)} damunt el
\VUgras{mur} \textsuperscript{(33)}, i les seves
\VUgras{roses} \textsuperscript{(39)}, que formen una tanca tan bonica. Contempla
els rics vestits de les senyores del seguici.

%%K t03-c1-15-1 p
15. --- Els nois del poble no fan gaire cas dels vestits. Es precipiten i s'empenyen per
aconseguir \VUgras{caramels} \textsuperscript{(55)} o
\VUgras{monedes} \textsuperscript{(52)}. Un d'ells plora
\textsuperscript{(40)}.

%%K t03-c1-16-1 p
16. --- Un altre ha trepitjat la pota del \VUgras{gos} \textsuperscript{(42)}, que fuig
xisclant i fa fugir la \VUgras{gallina} \textsuperscript{(43)} i els seus
\VUgras{pollets} \textsuperscript{(44)}.

%%K t03-tit-6 sub
\VUsaut{5.0mm}
\VUpk{10.2pt}{40}{El seguici del bateig.}
\VUsaut{4.0mm}

%%K t03-c1-17-1 p
17. --- Davant del seguici camina la \VUgras{dida} \textsuperscript{(45)}. Porta una bonica
\VUgras{còfia} \textsuperscript{(46)} amb cintes llargues. Porta als braços el \VUgras{nadó}
\textsuperscript{(47)}, tot vestit de \VUgras{puntes} \textsuperscript{(48)}.

%%K t03-c1-18-1 p
18. --- Darrere de la dida vénen el \VUgras{padrí} \textsuperscript{(49)} i la
\VUgras{padrina} \textsuperscript{(53)}. Reparteixen els caramels i les monedes.
El padrí porta \VUgras{guants} \textsuperscript{(51)} i \VUgras{barret de copa}
\textsuperscript{(50)}. La padrina porta a la mà un bonic \VUgras{ram} \textsuperscript{(54)}.

%%K t03-c1-19-1 p
19. --- Darrere d'ells veiem l'\VUgras{oncle} \textsuperscript{(56)} i la
\VUgras{tia} \textsuperscript{(58)} del nen. La tia porta un elegant
\VUgras{barret} \textsuperscript{(59)}, l'oncle una \VUgras{levita}
\textsuperscript{(57)} negra i una armilla blanca. Vénen després el \VUgras{cosí}
\textsuperscript{(60)} i la \VUgras{cosina} \textsuperscript{(61)}. Aquesta porta de la mà la
\VUgras{germana} \textsuperscript{(62)} del nadó, la petita Berta.

%%K t03-c1-20-1 p
20. --- L'altre \VUgras{germà} \textsuperscript{(63)} de Berta camina darrere. Dóna la mà a
una \VUgras{parenta} \textsuperscript{(64)}. Els \VUgras{amics} \textsuperscript{(65)} de la
família vénen després.

%%K t03-tit-7 sub
\VUsaut{4.6mm}
\VUpk{10.2pt}{40}{Els qui miren.}
\VUsaut{3.4mm}

%%K t03-c1-21-1 p
21. --- Prop del seguici, un \VUgras{captaire} \textsuperscript{(66)} espera recolzat a la
seva \VUgras{crossa} \textsuperscript{(67)}. Demana almoina.

%%K t03-c1-22-1 p
22. --- A l'altre costat hi ha un nen molt \VUgras{cortès} \textsuperscript{(68)}. Saluda i
dóna els \VUgras{«bon dia»} a les persones que coneix.

%%K t03-c1-23-1 p
23. --- Els del poble han sortit de casa seva per veure passar el seguici del bateig. Veiem un
\VUgras{pagès} \textsuperscript{(69)} amb la seva \VUgras{gorra de cotó} i al seu costat una
\VUgras{pagesa} \textsuperscript{(70)}. Després una
\VUgras{vella} \textsuperscript{(71)} recolzada al seu \VUgras{bastó}
\textsuperscript{(72)}. Finalment, el \VUgras{moliner} \textsuperscript{(73)} amb el seu ample
\VUgras{barret de feltre} \textsuperscript{(74)}, i una pagesa jove calçada amb
\VUgras{esclops} \textsuperscript{(75)}, que ensenya de caminar al seu nen.

%%K t03-c1-24-1 p
24. --- És una escena molt divertida.
\VUsaut{4.0mm}

%%K t03-tit-8 sub
\VUcentre{12.0pt}{0}{\textit{Segona escena.}}
\VUsaut{3.0mm}

%%K t03-tit-9 sub
\VUpk{10.2pt}{40}{Una festa pública.}
\VUsaut{2.4mm}

%%K t03-tit-10 sub
\VUpk{10.2pt}{40}{Jocs nàutics i regates.}
\VUsaut{3.0mm}

%%K t03-c2-01-1 p
1. --- Ha arribat l'\VUgras{estiu}. El sol ardent del mes de \VUgras{juny} (juliol o
\VUgras{agost}) convida els joves a banyar-se i a sortir en barca.

%%K t03-c2-02-1 p
2. --- A la riba dreta del riu, els remers es despullen per prendre part en els jocs nàutics i
en les regates.

%%K t03-c2-03-1 p
3. --- Un d'ells es treu les \VUgras{sabates} \textsuperscript{(1)}; les descorda (les
descargola). Els seus dos companys ja són a punt. No porten més que el seu
\VUgras{jersei} \textsuperscript{(2)} i el seu \VUgras{banyador}
\textsuperscript{(3)}. Xerren mentre esperen els seus amics. Parlen de la fira i de la festa
que se celebra a l'altra riba.

%%K t03-c2-04-1 p
4. --- A terra, al costat d'ells, jeu la \VUgras{roba} dels seus companys: un
\VUgras{parell} de \VUgras{botins} \textsuperscript{(4)}, \VUgras{sabates}
\textsuperscript{(5)}, \VUgras{mitjons} \textsuperscript{(6)}, barrets de
\VUgras{feltre} \textsuperscript{(7)} i de \VUgras{palla} \textsuperscript{(8)} i
una \VUgras{corbata} \textsuperscript{(9)}.

%%K t03-c2-05-1 p
5. --- Un dels joves ha plegat amb cura la seva \VUgras{jaqueta} \textsuperscript{(10)}. La
posa sobre una tovallola, en la qual el seu veí embolicarà de seguida els dos
\VUgras{vestits}.

%%K t03-c2-06-1 p
6. --- Un sisè noi va amb retard. Encara porta la seva \VUgras{gorra} \textsuperscript{(11)}
al cap. No s'ha tret ni la \VUgras{camisa} \textsuperscript{(12)}, ni el \VUgras{coll}
\textsuperscript{(13)}, ni els
\VUgras{punys} \textsuperscript{(14)}. Té a la mà la seva \VUgras{armilla}
\textsuperscript{(17)} i porta encara els seus \VUgras{pantalons} \textsuperscript{(16)} i els
seus \VUgras{elàstics} \textsuperscript{(15)}. Segur que no estarà a punt per a la pròxima
cursa.

%%K t03-c2-07-1 p
7. --- Prop dels joves, un caminet vorejat de \VUgras{cards} \textsuperscript{(18)} baixa a la
riba del riu, on l'oncle Jaume prepara, entre les \VUgras{canyes} \textsuperscript{(19)}, els
\VUgras{focs artificials} \textsuperscript{(20)} que es dispararan al vespre.

%%K t03-tit-11 sub
\VUsaut{4.4mm}
\VUpk{10.2pt}{40}{La cursa.}
\VUsaut{3.4mm}

%%K t03-c2-08-1 p
8. --- Al riu, algunes senyores que s'aixopluguen sota les seves ombrel·les passegen en
\VUgras{barca} \textsuperscript{(21)}. Segueixen la cursa, que és molt interessant. A cada
\VUgras{iol} \textsuperscript{(22)}, quatre
\VUgras{remers} \textsuperscript{(24)} remen amb totes les seves forces, mentre
el \VUgras{timoner} \textsuperscript{(26)} sosté el \VUgras{timó} \textsuperscript{(25)} i
dirigeix la fràgil embarcació. Els \VUgras{rems} \textsuperscript{(23)} colpegen l'aigua a
compàs, i els lleugers bots naveguen a una velocitat vertiginosa.

%%K t03-c2-09-1 p
9. --- Uns quants joves es disputen, vora el pilar del pont, el premi de la
\VUgras{cucanya horitzontal} \textsuperscript{(28)}. Un d'ells avança amb cautela
pel pal flexible i lliscós per atènyer la petita \VUgras{bandera} \textsuperscript{(29)}
subjectada a l'extrem. El seu dissortat
\VUgras{competidor} \textsuperscript{(30)} ha caigut a l'aigua. Neda per arribar
a la riba.

%%K t03-c2-10-1 p
10. --- Uns nois més grans \VUgras{justen en barca} \textsuperscript{(32)} vora el
\VUgras{moll} \textsuperscript{(33)}. A tots aquests jocs hi assisteix un
\VUgras{públic} \textsuperscript{(31)} nombrós, recolzat al \VUgras{ampit} del
\VUgras{pont} \textsuperscript{(27)} i atapeït a la \VUgras{riba}. Alguns
espectadors, però, es giren per seguir el joc de la \VUgras{cucanya} \textsuperscript{(45)} o
per admirar el \VUgras{seguici municipal} que acut al
\VUgras{repartiment} dels premis.

%%K t03-c2-11-1 p
11. --- Primer van uns \VUgras{trinxeraires} \textsuperscript{(35)} al davant dels
\VUgras{músics} \textsuperscript{(36)}; després vénen l'\VUgras{alcalde}
\textsuperscript{(37)}, els tinents d'alcalde i l'\VUgras{ajuntament} de la petita ciutat.
Tanquen la marxa els \VUgras{bombers} \textsuperscript{(38)}.

%%K t03-tit-12 sub
\VUsaut{5.0mm}
\VUpk{10.2pt}{40}{La fira.}
\VUsaut{3.6mm}

%%K t03-c2-12-1 p
12. --- Hi ha molta gent al \VUgras{camp de la fira} \textsuperscript{(39)}. Es ve a veure les
diferents \VUgras{barraques}: els \VUgras{cavallets de fusta} \textsuperscript{(40)}, el
\VUgras{circ} \textsuperscript{(41)}, on la funció ja ha començat; el \VUgras{ball públic}
\textsuperscript{(42)}, on els joves i les joves de la ciutat gaudeixen del plaer de
\VUgras{ballar}.

%%K t03-c2-13-1 p
13. --- Al fons veiem la \VUgras{plaça de braus} \textsuperscript{(44)}, on el valent
\VUgras{matador} espanyol lluita contra el \VUgras{brau} \textsuperscript{(45)} furiós;
després la \VUgras{muntanya russa} \textsuperscript{(49)} i la \VUgras{barraca de tir}
\textsuperscript{(46)}, on s'exerciten a manejar les \VUgras{armes de foc} i a tirar al
\VUgras{blanc}.

%%K t03-c2-14-1 p
14. --- Prop d'allà hi ha el \VUgras{teatre de fira} \textsuperscript{(47)}, en què es
representen la \VUgras{comèdia} i el \VUgras{drama}. Molt a prop hi ha la barraca del
\VUgras{gegant} \textsuperscript{(48)} i del \VUgras{nan}. L'\VUgras{estatura} del primer és
de dos metres quinze; el segon amb prou feines és tan alt com un nen.

%%K t03-c2-15-1 p
15. --- I al final de tot, la gran \VUgras{col·lecció de feres} \textsuperscript{(50)} amb els
seus \VUgras{animals salvatges}, el seu
\VUgras{tigre} \textsuperscript{(51)} de poderoses \VUgras{urpes}, el seu
\VUgras{lleó} \textsuperscript{(52)} d'espessa \VUgras{cabellera}, les seves dues
\VUgras{girafes} \textsuperscript{(53)} i el seu \VUgras{elefant}
\textsuperscript{(54)}, que fa servir amb tanta traça la seva \VUgras{trompa}.

%%K t03-c2-16-1 p
16. --- El \VUgras{domador} \textsuperscript{(55)} que ensinistra aquests animals és a la
porta de la barraca.
\VUsaut{6.0mm}
\VUfilet{22mm}
